%\thispagestyle{empty}
\section*{Abkürzungsverzeichnis}
\addcontentsline{toc}{section}{Abkürzungsverzeichnis}

%% Alle Erweiterungen sind aus dem Abkürzungsverzeichnis der Vorarbeiten
%% (Friesen, MA_Bericht_Friesen.pdf, S. VIII-IX) übernommen, um im Bericht
%% konsistente Bezeichnungen zu verwenden. Bei Bedarf bitte um weitere
%% Einträge (z.B. KCB, KLI, KPC, WVS) ergänzen -- diese habe ich hier
%% weggelassen, da sie im aktuellen Text nicht vorkommen.

\begin{acronym}
    %C
    \acro{circ}[CIRC]{kreisförmige Bewegung}
    %E
    \acro{ea}[E/A]{Ein-/Ausgabe}
    \acro{ethercat}[EtherCAT]{Ethernet for Control Automation Technology}
    \acro{ethip}[Ethernet/IP]{EtherNet Industrial Protocol}
    %G
    \acro{gui}[GUI]{Graphical User Interface}
    %K
    \acro{krl}[KRL]{KUKA Robot Language}
    \acro{ksi}[KSI]{KUKA Service Interface}
    \acro{kvp}[KVP]{KukaVarProxy}
    %L
    \acro{lin}[LIN]{lineare Bewegung}
    %O
    \acro{opc}[OPC]{Open Platform Communications}
    \acro{opcua}[OPC~UA]{Open Platform Communications Unified Architecture}
    %P
    \acro{ptp}[PTP]{Point-to-Point}
    %Q
    %% QUARC ist ein Produktname (MathWorks/Quanser), kein Akronym im
    %% eigentlichen Sinn -- daher hier bewusst nicht aufgenommen.
    %R
    \acro{rsi}[RSI]{Robot Sensor Interface}
    %S
    \acro{spl}[SPL]{Spline-Bewegung}
    \acro{sps}[SPS]{Speicherprogrammierbare Steuerung}
    %T
    \acro{tcp}[TCP]{Tool Center Point}
    \acro{tcpip}[TCP/IP]{Transmission Control Protocol / Internet Protocol}
    %U
    \acro{udp}[UDP]{User Datagram Protocol}
\end{acronym}
